% ══════════════════════════════════════════════════════════════════════════
\section{What this paper borrows}
\label{sec:model}
% ══════════════════════════════════════════════════════════════════════════

Everything in this section is developed in the companion~\cite{deadly-race} and
restated here only so that this paper can be read alone. Notation: $\Victim$ is
the holder, $\Attacker$ the coercer, $\Kset$ the \emph{spend authority} --- the
set of states from which the stash can be spent, a single key being one case ---
and $t$ the coercion-time budget.

\begin{assumption}[The Wrench--Dolev--Yao (WDY) adversary]\label{ass:power}
$\Attacker$ knows the protocol and the design landscape (Kerckhoffs), can seize
every device co-located with $\Victim$, and can coerce $\Victim$ for a bounded
time $t$. He is \emph{priced}: he holds a value $W$ for the stash and a cost $c$
for a homicide, and he acts to maximise expected return. The construction is
\emph{a fortiori} --- he is granted strictly more than a real coercer, so a bar
cleared against him is cleared against the real one.
\end{assumption}

\begin{assumption}[Non-Negative-Proof Principle, NNPP]\label{lem:nnpl}
$\Victim$ cannot prove the non-existence of an undisclosed backup. Non-possession
is a negative existential over everything in $\Victim$'s reach, the last site
being whether $\Victim$ has forgotten or never memorised. $\Attacker$ therefore
cannot be brought to certainty by any disclosure $\Victim$ makes.
\end{assumption}

\noindent Two consequences of NNPP carry the weight here. First, cooperation
cannot dispel the incentive: no act of disclosure closes the question, so
compliance buys no guaranteed release~\cite[Cor.~1]{deadly-race}. We refer to
this as the \textbf{disclosure corollary}.\label{cor:disclosure} Second, a denial
is a claim $\Attacker$ has both an incentive to disbelieve and no means of
settling --- which is the hinge of \S\ref{sec:spiral}.

\begin{definition}[Naive coercion resistance]\label{def:naive-cr}
A protocol is \textbf{naively coercion-resistant} if it \emph{may} resist a
\emph{sub-Kerckhoffs} attacker --- one with the physical capabilities but lacking
knowledge of the design --- because denial, decoy or obscurity may work when the
attacker does not know $\Victim$'s setup. Both hedges are load-bearing: the
guarantee is contingent on the attacker's \emph{ignorance}, not on any structural
property, and so offers none against the full WDY adversary. This is the tier the
whole obscurity class occupies.
\end{definition}

\noindent Where the companion's stronger tiers are needed we name them in prose:
a protocol is \emph{weakly} resistant if a single encounter leaves a
feasible-but-unfinished path and so a positive-pivotality state, and
\emph{strongly} resistant if no reachable state has positive
pivotality~\cite[\S3]{deadly-race}.

\begin{lemma}[Deadly Race, after~{\cite[Lem.~1]{deadly-race}}]\label{lem:drl}
If, on release, $\Attacker$ holds a feasible path to $\Kset$ while $\Victim$
retains one too, then $\Victim$ is \emph{pivotal}: removing him raises
$\Attacker$'s win probability by $\pi > 0$, and the removal pays whenever
$(1-p_A)\,W > c$. The \textbf{No-Gray-Area Criterion} is the resulting design
bar --- admit no such state.\label{prin:nga}
\end{lemma}

\noindent Finally, the companion's \emph{Four Properties} are used here only to
locate the obscurity class. They are \textbf{(\ref{p:know})} knowledge-based
authentication, \textbf{(\ref{p:indiv})} individual custody,
\textbf{(\ref{p:nonob})} non-obscurity --- which \emph{is} the Kerckhoffs premise
--- and \textbf{(\ref{p:coerc})} coercion resistance:
\begin{enumerate}[nosep,leftmargin=2.4em,label=\textbf{(\arabic*)},ref=\arabic*]
  \item\label{p:know} knowledge-based authentication (deviceless);
  \item\label{p:indiv} individual custody;
  \item\label{p:nonob} non-obscurity;
  \item\label{p:coerc} coercion resistance.
\end{enumerate}
Clearing the bar while holding all four forces the \emph{tacit}
route~\cite[Cor.~3]{deadly-race}; the alternative that keeps a winning racer
beyond the attacker's reach we call the \textbf{remote-delegate}
route.\label{r:remotewinner}
